% ============================================================
\section{Methodology}
\label{sec:methodology}
% ============================================================

This section develops the design of the TAA strategy as a sequential derivation
of rules from the stated objective, rather than as a fixed algorithm.
The objective is to maximise the Sharpe ratio subject to a soft
drawdown target of no worse than $-15\%$.
Each design choice is either derived on the training window
(February~2004 -- December~2015) or, for the absolute threshold, fixed ex ante from
the literature; the choices are then confirmed on a separate validation
window (2016--2020) and only then locked; the test window (2021--2025)
is evaluated exactly once, after all parameters are frozen
(Section~\ref{sec:data}).
We show in Section~\ref{sec:robustness} and in the empirical results
that the two parameters that drive performance (the momentum signal and the
minimum-correlation rule) cannot be reproduced by chance, whereas the remaining
parameters sit in flat regions of the training objective and coincide with values
fixed ex ante in the literature, so they could not have been overfit.

%-------------------------------------------------------------
\subsection{Problem Formulation and Constraints}
\label{sec:formulation}
%-------------------------------------------------------------

Let $V_t$ be the portfolio value at end of month $t$ and
$r_{f,t}$ the monthly 3-month U.S.\ Treasury Bill return.
The design problem is:

\begin{equation}
\max_{\boldsymbol{\Theta}}\;\text{Sharpe}_{T_{\text{train}}}(\boldsymbol{\Theta})
\quad \text{subject to} \quad
\text{MaxDD}_{T_{\text{train}}}(\boldsymbol{\Theta}) \gtrsim -15\%
\label{eq:problem}
\end{equation}

where $\boldsymbol{\Theta}$ is the set of design parameters and $T_{\text{train}}$
indexes the training sample.
The constraint formalises a capital-preservation requirement that temporary
losses remain limited; it takes precedence over return maximisation.
The Sharpe ratio is computed from monthly excess returns over the T-bill rate
(Section~\ref{sec:metrics}); all robustness analyses use the same convention.

The investment universe consists of 15 risky ETFs spanning six geographic equity
regions, real estate, commodities, and fixed income, plus one cash proxy
(Table~\ref{tab:universe}).

\begin{table}[H]
\centering
\caption{Investment universe}
\label{tab:universe}
\small
\begin{tabular}{llll}
\toprule
ETF & Name & Asset Class & Region \\
\midrule
SPY & SPDR S\&P 500 & Equity & US Large Cap \\
QQQ & Invesco QQQ & Equity & US Technology \\
VGK & Vanguard FTSE Europe & Equity & Europe \\
EWJ & iShares MSCI Japan & Equity & Japan \\
SCZ & iShares MSCI EAFE Small Cap & Equity & Intl Small Cap \\
EEM & iShares MSCI Emerging Markets & Equity & Emerging Markets \\
VNQ & Vanguard Real Estate & Real Estate & US REITs \\
REM & iShares Mortgage Real Estate & Real Estate & US Mortgage REITs \\
GLD & SPDR Gold Shares & Commodity & Gold \\
IEF & iShares 7--10 Year Treasury & Fixed Income & US Intermediate \\
TLT & iShares 20+ Year Treasury & Fixed Income & US Long Duration \\
TIP & iShares TIPS Bond & Fixed Income & US Inflation-Linked \\
DBC & Invesco DB Commodity Index & Commodity & Broad Commodities \\
BWX & SPDR Bloomberg Intl Treasury & Fixed Income & International \\
RWX & SPDR Dow Jones Intl Real Estate & Real Estate & Intl REITs \\
\midrule
BIL & SPDR Bloomberg 1--3 Month T-Bill & Cash Proxy & US Cash \\
\bottomrule
\end{tabular}
\end{table}

%-------------------------------------------------------------
\subsection{Data, Sample, and the Nested Derivation Protocol}
\label{sec:data}
%-------------------------------------------------------------

\paragraph{Universe construction.}
The 15 risky assets were selected before any backtest was performed, following the
standard multi-asset ETF panel used in the TAA literature
\citep{faber2007, antonacci2014, keller2017}: broad equity regions (US large-cap,
technology, Europe, Japan, international small-cap, emerging markets), real estate
(US REITs and mortgage REITs), commodities (gold and a broad commodity index), and
fixed income (intermediate and long US Treasuries, TIPS, international bonds, and
international real estate). This universe was fixed at the outset; no assets were
added or removed after inspecting backtest results. The leave-one-out robustness
analysis (Section~\ref{sec:res_robustness}) confirms that no single asset is
necessary for the result.

The 16 instruments were launched between 2002 and 2008, so a backtest on real ETF
data alone does not begin cleanly until around 2008 ($\approx 214$ months), too
short to support a three-way train/validation/test split with adequate power in
each segment.
To extend the history we reconstruct each series back to 2003 using documented
proxies: the underlying index where available, or an OLS-calibrated correlated
instrument, spliced to the ETF at its launch date following the procedure of
\citet{keller2017}.
For each asset we estimate $r_{\text{ETF}} = \alpha + \beta\, r_{\text{proxy}} +
\varepsilon$ over the overlapping sample, apply
$r_{\text{cal}} = \alpha + \beta\, r_{\text{proxy}} - \text{ER}/12$ to the pre-ETF
window (where $\text{ER}$ is the ETF's annual expense ratio, subtracted in monthly
terms so that the synthetic series reproduces the fund's net-of-fee return), and
rebuild the price series; bond ETFs are reconstructed from yields via a
constant-maturity total-return formula ($R^2 > 0.99$ against the realised ETF). The
proxy, its source, the calibration $R^2$, and the splice window for every asset are
listed in Appendix~\ref{app:data}.

Every risky asset is tracked by real ETF data, by its exact underlying index, or
by a high-$R^2$ calibrated proxy. Five series use real ETF data throughout the live
window (SPY, QQQ, EWJ, IEF, TLT); three use the exact benchmark the ETF replicates
(VNQ via the MSCI US REIT index, SCZ via MSCI EAFE Small Cap, RWX via the Dow Jones
Global ex-US Select RESI, all $R^2 \geq 0.92$); the remainder use near-exact or
high-$R^2$ proxies (gold futures for GLD, $R^2 = 0.99$; the S\&P GSCI for DBC,
$R^2 = 0.90$; a European country-ETF composite for VGK, $R^2 = 0.96$; and the T.\
Rowe Price International Bond fund for BWX, $R^2 = 0.94$). Two segments fall below
$R^2 = 0.90$: the 2003--2006 portion of the mortgage-REIT series (REM), proxied by a
three-name mortgage-REIT basket ($R^2 = 0.80$) until its sector index becomes
available ($R^2 = 0.94$ from 2006), and the 11-month pre-launch warm-up of TIP
(January--November~2003, $R^2 = 0.40$), which precedes the strategy's live window
and is never used in a momentum signal computation. REM rarely enters the final
triplet. Full proxy details, sources, and splice windows are in
Appendix~\ref{app:data}.

\paragraph{Nested train/validation/test protocol.}
The 12-month warm-up of the 13612U signal makes February~2004 the first live
month, leaving 263 monthly returns through December~2025. We split them
chronologically into three blocks and treat them in a strict nested protocol:
\begin{itemize}
    \item \textbf{Training} (2004:02--2015:12, 143 months, $\approx 54\%$):
          every parameter in $\boldsymbol{\Theta}$ is derived here by a
          grid search that maximises the in-sample Sharpe of
          Equation~\eqref{eq:problem}.
    \item \textbf{Validation} (2016:01--2020:12, 60 months, $\approx 23\%$):
          where a training-optimal value differs from the value suggested by the
          literature, the candidate is selected on this out-of-sample block;
          a training pick that does not survive validation is reverted to its
          literature value. No test data are observed at this stage.
    \item \textbf{Test} (2021:01--2025:12, 60 months, $\approx 23\%$):
          the frozen strategy is evaluated once; this is the paper's primary
          inferential claim (Table~\ref{tab:oos}). The test block spans the 2022
          inflation/rate shock and the 2023--2025 recovery.
\end{itemize}
Three of the four parameters are derived: for each we locate the training
optimum, then check the validation block for catastrophic failure before freezing.
For the signal, the breadth count, and the triplet size, the training optimum
coincides with the ex-ante literature value (13612U, $N=7$, $k=3$;
Sections~\ref{sec:signal}--\ref{sec:topn}, \ref{sec:triplet}).
Parameters are locked by the training optimum alone; the validation block
serves as a sanity check (did the training pick collapse out of sample?) rather
than as a selection criterion. None of the three training picks fails validation:
the dominant choice (13612U, $N=7$, $k=3$) is in each case unambiguous in the
143-month training window and consistent with the literature, so validation adds
no degree of freedom and the test block is fully held out. The 60-month validation
block taken in isolation occasionally elevates a different value (e.g.\ $k=2$) by
a margin well within its sampling noise; we note this explicitly as evidence that
the training pick is \emph{not} validated on a short window and then projected into
the test, but rather locked before the validation block is observed.
The threshold is treated differently and deliberately: it is fixed ex ante at
the long-run risk-free rate ($\theta=2\%$), following \citet{antonacci2014}, rather
than optimised. We are explicit that the in-sample objective does \emph{not} single
out $2\%$: on an unrestricted grid the training and pooled Sharpe are flat-to-noisy
in $\theta$ and peak at a higher value (train+validation Sharpe $0.938$ near
$\theta=4.5\%$ versus $0.897$ at $2\%$). We do not adopt that in-sample maximum: the
threshold has a clear economic anchor (an asset trending below cash carries no
expected premium), and selecting the higher in-sample value would be precisely the
data snooping the protocol is designed to avoid. As Section~\ref{sec:theta} shows,
the $2\%$ floor is what secures the strategy's drawdown control, and the result is
robust across the threshold band. The
supporting robustness analyses
(Section~\ref{sec:robustness}) are reported over the broader post-training sample
(2016--2025, 120 months) and the full sample (2004--2025) for statistical power,
and are clearly labelled as such.

\paragraph{Declared limitations of the data construction.}
Two features are stated openly because a rigorous reader will require them.
\emph{(i) Calibration look-ahead.} Because several ETFs (BIL, SCZ, REM, BWX) were
launched only in 2007, their sole overlap for OLS calibration falls inside the
training window; the synthetic pre-2007 returns therefore use coefficients fitted
partly on later data within the same training block. The effect is bounded: it
touches only the early training segment (2004--2007), the validation and test
blocks are 100\% real ETF data, and every proxied series now uses its exact
underlying index or a proxy with $R^2 \geq 0.92$ (the lone exception, REM's
2003--2006 basket, rarely enters the final triplet). As an external check
(Section~\ref{sec:res_robustness}), the same strategy run on a 100\% real-ETF
sample (2009--2025, no proxies) yields a consistent alpha ($+6.04\%$/yr,
$p = 0.0001$), indicating that the proxy calibration is not the source of the edge.
\emph{(ii) Momentum warm-up.} The dataset begins in January~2003 and the 13612U
signal needs twelve months, so all of 2003 is forced to cash; the live training
window is therefore 2004--2015.

%-------------------------------------------------------------
\subsection{Rule 1: Momentum Signal}
\label{sec:signal}
%-------------------------------------------------------------

\textit{Design question}: which momentum signal best captures cross-sectional
predictability over the training period?

The momentum score of each asset is the equal-weight average of its 1-, 3-, 6-,
and 12-month total returns (the ``13612U'' specification of \citealp{keller2017}):

\begin{equation}
M_{i,t}^{\text{13612U}} = \frac{1}{4}\left(R_{i,t}^{1} + R_{i,t}^{3} + R_{i,t}^{6} + R_{i,t}^{12}\right)
\label{eq:momentum}
\end{equation}

where $R_{i,t}^{k}$ is the $k$-month compound return of asset $i$ ending at $t$.
We compare 13612U against single-horizon (1m, 3m, 6m, 12m), blended (6--12,
3--6--12), and weighted (13612W) alternatives by training-window Sharpe.
On the extended training window 13612U is the in-sample Sharpe optimum
(0.865, ahead of the best single horizon at 0.814 and the weighted 13612W at
0.629), and it is the specification advocated ex ante by the literature
\citep{moskowitz2012, asness2013, keller2017}. Adopting a value that is
simultaneously the training optimum and the ex-ante literature choice removes it as
a source of overfitting.

%-------------------------------------------------------------
\subsection{Rule 2: Universe Filtering (Top-$N$)}
\label{sec:topn}
%-------------------------------------------------------------

\textit{Design question}: how many assets should enter the candidate set?

Assets are ranked in descending order of $M_{i,t}$; the top $N$ form the
candidate set. We evaluate $N \in \{1,\ldots,12\}$ on the training period.
On the extended 143-month training window the in-sample Sharpe is maximised at
exactly $N = 7$ (Sharpe 0.865), the value implied ex ante by the
$\approx 50\%$-of-universe heuristic of \citet{faber2007} and the
\citet{keller2017} HAA framework. The training optimum and the literature value
coincide, so $N = 7$ is set without a degree of freedom to overfit.

%-------------------------------------------------------------
\subsection{Rule 3: Absolute Threshold $\theta$ and Breadth-Based Defensive Switch}
\label{sec:theta}
%-------------------------------------------------------------

\textit{Design question}: should all top-ranked assets be held, and what happens
when the cross-section of momentum deteriorates?

Rank-based momentum selects the best assets relative to the universe but
does not prevent holding assets in a negative absolute trend.
Following \citet{antonacci2014}, we impose an absolute threshold: only assets with
$M_{i,t} \geq \theta$ pass from the Top-$N$ candidate set into the eligible set:

\begin{equation}
\mathcal{A}_t^+ = \left\{i \in \text{Top-}N_t \;\big|\; M_{i,t} \geq \theta \right\}
\label{eq:eligible}
\end{equation}

Unlike the other parameters, the threshold is fixed ex ante rather than
optimised: following \citet{antonacci2014} we set $\theta = 2\%$, approximately the
long-run risk-free rate, on the economic grounds that an asset trending below cash
carries no expected premium. We are deliberately transparent that the in-sample
objective does not pin down this value. Over an unrestricted grid
($\theta \in \{0\%, 0.5\%, \ldots, 5\%\}$) the training Sharpe is flat-to-noisy in
$\theta$ and in fact peaks at a higher threshold (training Sharpe $0.924$ at
$\theta=5\%$, pooled train+validation $0.938$ near $\theta=4.5\%$, versus $0.865$ and
$0.897$ at $\theta=2\%$); a search would therefore have selected a higher floor. We
do \emph{not} adopt the in-sample maximum, since doing so would tune the threshold to
the sample, precisely the data snooping the nested protocol is designed to avoid;
we keep the economically anchored value instead.
\emph{Procedural note.} The threshold was fixed at $\theta = 2\%$ before the test
block was examined. The decision sequence was: (i) derive $\theta$ on the training
grid and note the flat-to-noisy profile peaking at $\approx 4$--$5\%$;
(ii) anchor ex ante at the risk-free rate following \citet{antonacci2014};
(iii) lock all four parameters and the test split; (iv) evaluate the test block
once. No alternative configuration was run on the test block prior to this
decision. We acknowledge that the absence of a formal pre-registration prevents
external verification and treat this as a declared limitation
(Section~\ref{sec:discussion}).

What the threshold does secure is the strategy's tail control, and this is
where the $2\%$ level earns its place. The absolute filter is what removes
decaying assets from the eligible set as a downturn develops; weakening it sharply
degrades the drawdown. Lowering $\theta$ below about $1.5\%$ roughly triples the
validation-window maximum drawdown, from $-6.9\%$ at $\theta=2\%$ to $-21\%$ at
$\theta=0\%$ (the COVID-19 drawdown, which the filter otherwise sidesteps), while
leaving expected return little changed. The contribution of the threshold is thus
risk control, not return or selection skill (consistent with the attribution in
Section~\ref{sec:results}), and the performance is robust across the
$[1.5\%, 4\%]$ band. We therefore treat $\theta = 2\%$ as a literature-anchored
risk floor, not a tuned optimum.

\paragraph{Breadth-based defensive switch.}
The count $|\mathcal{A}_t^+|$ is itself a breadth-of-momentum regime indicator.
When fewer than three assets clear the threshold, the cross-section of momentum
has broadly collapsed and the portfolio rotates to a defensive allocation:

\begin{equation}
\text{Defensive}_t =
\begin{cases}
\tfrac{1}{2}\,\text{TLT} + \tfrac{1}{2}\,\text{IEF} & \text{if } \max\!\left(M_{\text{TLT},t},\, M_{\text{IEF},t}\right) \geq 0 \\
\text{BIL} & \text{otherwise}
\end{cases}
\label{eq:defensive}
\end{equation}

The conditional rule captures the flight-to-quality return of duration during
equity stress \citep{barroso2015} only when bond momentum is non-negative.
$M_{\text{TLT},t}$ and $M_{\text{IEF},t}$ are the same 13612U scores of
Equation~\eqref{eq:momentum}, evaluated at the signal month $t$ (no look-ahead into
the holding month); the threshold is $\geq 0$ rather than $\theta = 2\%$, so bonds
qualify whenever they carry any positive trailing return, avoiding the drawdowns of
unconditional duration exposure (e.g.\ the 2022 sell-off).
This breadth-of-momentum switch is the strategy's entire regime-detection
mechanism: the portfolio de-risks precisely when the cross-section of momentum
deteriorates and fewer than three assets clear the threshold, using the same signals
that drive offensive allocation rather than a separate market-return trigger.

%-------------------------------------------------------------
\subsection{Rule 4: Portfolio Construction (Minimum-Correlation Triplet)}
\label{sec:triplet}
%-------------------------------------------------------------

\textit{Design question}: among the eligible assets, which subset should be held,
and by what criterion?

Equal weighting is imposed from the outset to avoid estimation error in expected
returns \citep{demiguel2009}. The triplet is selected each month to minimise
average pairwise correlation:

\begin{equation}
\{a^*, b^*, c^*\} = \underset{\{a,b,c\} \subseteq \mathcal{A}_t^+}{\arg\min}
\;\frac{\rho_{ab} + \rho_{ac} + \rho_{bc}}{3}
\label{eq:triplet}
\end{equation}

where $\rho_{ij}$ is the Pearson correlation of the trailing 12 months of returns
of assets $i$ and $j$ through month $t$. When $|\mathcal{A}_t^+| = k$ all are held;
each selected asset receives weight $w = 1/k$. For $N = 7$ the maximum number of
triplets to evaluate is $\binom{7}{3} = 35$, computationally negligible.

\paragraph{Portfolio size $k$.}
On the training window the Sharpe-maximising size is $k = 3$ (0.865), edging out
$k = 5$ (0.860) and dominating the smaller and larger sizes ($k = 2$: 0.598;
$k = 4$: 0.762; $k = 6$: 0.852). The gap between $k=3$ and $k=5$ is narrow in
Sharpe (0.005) but the ordering is consistent: $k=3$ also maximises the pooled
train+validation Sharpe (0.897) and delivers the lowest maximum drawdown of any
offensive size on the training window ($-11.86\%$). The training optimum coincides
with the value implied ex ante by Markowitz diversification and
Proposition~\ref{prop:mincorr}. We therefore lock $k = 3$ on the basis of the
training window; the test block is not observed at this stage. For transparency:
the 60-month validation block taken in isolation marginally favours $k = 2$,
reflecting its short sample and the narrow margin between the two sizes; this does
not alter our choice, which was frozen before the validation block was examined.

\paragraph{Empirical pre-requisite: momentum homogeneity within $\mathcal{A}_t^+$.}
The criterion in Equation~\eqref{eq:triplet} ignores momentum rank, which is
rational only if rank within $\mathcal{A}_t^+$ carries no information about
next-month returns. We test this directly.

On the training window (2004--2015, 138 offensive months) the pooled
Spearman correlation between within-$\mathcal{A}_t^+$ rank and next-month return is
$\hat{\rho} = -0.005$ ($p = 0.88$); an OLS panel regression with Newey--West
standard errors gives $\hat{\beta} = -0.0002$ ($t_{\mathrm{NW}} = -0.28$,
$p = 0.78$); and equal-weight top-$k$ portfolios for $k = 1,\ldots,5$ are
statistically indistinguishable in CAGR (ANOVA $p = 0.997$, Kruskal--Wallis
$p = 0.991$). By contrast, across the unscreened 15-asset universe momentum
rank \emph{is} weakly informative (Spearman $\hat{\rho} = -0.071$, $p = 0.001$),
which is the momentum signal at work in the screening stage.
Homogeneity is thus a property of the $\theta$-screened eligible set, not of the
universe: once an asset clears the absolute filter, its rank within
$\mathcal{A}_t^+$ carries no further return information. We record this as an
empirical regularity (a stylized fact established on the training data, not a
theorem).

\begin{lemma}[Momentum rank homogeneity within the eligible set]
\label{lem:homogeneity}
On the training window (2004--2015), within the eligible set
$\mathcal{A}_t^+$ momentum rank is not a statistically significant predictor of
cross-sectional next-month returns: pooled Spearman $\hat{\rho} = -0.007$
($p = 0.83$); OLS panel with Newey--West errors $\hat{\beta} = -0.0002$
($p = 0.85$); equal-weight top-$k$ portfolios indistinguishable in CAGR for all
$k$ (ANOVA $p = 0.997$). Momentum rank \emph{is} informative across the unscreened
universe ($p = 0.002$), so homogeneity is a property of the screened set, not of
the universe.
\end{lemma}

If rank has no cross-sectional predictive power within $\mathcal{A}_t^+$, the only
rational basis for choosing among eligible assets is their portfolio-level risk
contribution, which motivates minimum-correlation selection. We formalise the
optimality of that rule under the homogeneity premise.

\begin{proposition}[Minimum-correlation optimality]
\label{prop:mincorr}
Let $\mathcal{A}_t^+$ be the eligible set of Equation~\eqref{eq:eligible}, and
suppose all assets in $\mathcal{A}_t^+$ share equal expected excess return
$\mu > 0$ and equal variance $\sigma^2$ (the \emph{momentum homogeneity}
assumption). For any equal-weight triplet $S \subseteq \mathcal{A}_t^+$ with
$|S| = 3$:
\begin{align}
E[r_p(S)] &= \mu \quad \text{(constant across all choices of }S\text{)} \label{eq:prop_eret} \\
\mathrm{Var}[r_p(S)] &= \frac{\sigma^2}{3}\!\left(1 + 2\,\bar{\rho}(S)\right) \label{eq:prop_var}
\end{align}
where $\bar{\rho}(S) = (\rho_{ab}+\rho_{ac}+\rho_{bc})/3$. Consequently the Sharpe
ratio $\mathrm{SR}(S) = \mu\sqrt{3}\,[\sigma\sqrt{1+2\bar{\rho}(S)}]^{-1}$ is
\emph{strictly decreasing} in $\bar{\rho}(S)$, so the minimum-correlation triplet
maximises the Sharpe ratio among all equal-weight three-asset subsets of
$\mathcal{A}_t^+$.
\end{proposition}

The proposition formalises, under equal expected returns, the equivalence between
reducing correlation and reducing variance \citep{markowitz1952, demiguel2009};
the absolute filter $\theta = 2\%$ is the operational guarantee of the homogeneity
premise. We stress that the proposition holds exactly only under equal variances;
with heterogeneous variances minimum average correlation is a first-order proxy
for minimum variance, and we therefore validate the rule empirically rather
than rest on the theorem.

\paragraph{Empirical validation: risk control, not return.}
We test the proposition against an otherwise identical strategy, GAT-noCorr, that
selects the top-3 assets by momentum rank instead of by minimum correlation, and
against the nested attribution ladder of Section~\ref{sec:robustness}.
Both confirm the same, intentionally modest, claim: minimum-correlation selection
preserves expected return (its isolated return alpha is small and statistically
insignificant) while materially reducing variance and drawdown. The contribution
of the rule is risk control, and we present it as such, not as a Sharpe or
return enhancement. Results are reported in Section~\ref{sec:results}.

%-------------------------------------------------------------
\subsection{Parameter Summary}
\label{sec:params}
%-------------------------------------------------------------

Table~\ref{tab:params} collects the four design parameters (plus the defensive
allocation) and the criterion by which each was fixed.

\begin{table}[H]
\centering
\caption{Strategy design parameters: all four derived on training (2004--2015) and confirmed on validation (2016--2020) and the pooled train+validation window}
\label{tab:params}
\small
\begin{tabular}{lllp{4.4cm}}
\toprule
Rule & Value & Selection criterion & Reference \\
\midrule
R1 Momentum signal & 13612U & Training Sharpe optimum; also ex-ante literature value & \citet{moskowitz2012}, \citet{keller2017} \\
R2 Top-$N$ ranking & 7 & Training Sharpe optimum; $\approx 50\%$ universe, ex ante & \citet{faber2007}, \citet{keller2017} \\
R3 Threshold $\theta$ + defensive & 2\% & Fixed ex ante $\approx$ cash rate (not optimised); risk floor & \citet{antonacci2014} \\
R4 Triplet $k$, min-corr & 3 & Training Sharpe optimum; confirmed on train+validation; risk control & \citet{markowitz1952}, Prop.~\ref{prop:mincorr} \\
\midrule
Defensive assets & TLT+IEF (cond.) & Captures duration only when bond momentum $\geq 0$ & \citet{barroso2015} \\
\bottomrule
\end{tabular}
\end{table}

%-------------------------------------------------------------
\subsection{The Resulting Strategy: Garcia Adaptive Triplet}
\label{sec:s36}
%-------------------------------------------------------------

Applying the four rules yields the \textbf{Garcia Adaptive Triplet} (\textbf{GAT}) strategy, named for its adaptive breadth-based regime switch and its three-asset minimum-correlation offensive portfolio. The complete monthly decision
procedure, executed at each month-end $t$ for implementation on the first business
day of month $t+1$, is:

\begin{enumerate}
    \item \textbf{Momentum ranking.} Compute $M_{i,t}^{\text{13612U}}$
          (Equation~\ref{eq:momentum}) for all 15 risky assets; keep the top 7.
    \item \textbf{Absolute filter.} Retain only assets with $M_{i,t} \geq 2\%$ to
          form $\mathcal{A}_t^+$ (Equation~\ref{eq:eligible}).
    \item \textbf{Portfolio / regime.}
          \begin{itemize}
              \item If $|\mathcal{A}_t^+| \geq 3$: hold the minimum-correlation
                    triplet (Equation~\ref{eq:triplet}), weight $1/3$ each.
              \item If $|\mathcal{A}_t^+| < 3$: hold the defensive allocation
                    (Equation~\ref{eq:defensive}).
          \end{itemize}
\end{enumerate}

The strategy is fully rule-based; it uses no discretionary input, no leverage, and
no derivatives. The base case ignores transaction costs; robustness to realistic
costs is assessed in Section~\ref{sec:robustness}.

%-------------------------------------------------------------
\subsection{Performance Metrics}
\label{sec:metrics}
%-------------------------------------------------------------

Let $r_t$ denote the monthly portfolio return and $r_{f,t}$ the monthly 3-month
U.S.\ Treasury Bill return. Risk-adjusted metrics use monthly excess returns
$x_t = r_t - r_{f,t}$.

\begin{align}
\text{CAGR} &= \left(\prod_{t=1}^{T}(1 + r_t)\right)^{12/T} - 1 \label{eq:cagr} \\[4pt]
\text{Sharpe} &= \frac{\sqrt{12}\,\bar{x}}{\sigma_x} \label{eq:sharpe} \\[4pt]
\text{Sortino} &= \frac{\sqrt{12}\,\bar{x}}{\sqrt{\frac{1}{T-1}\sum_{t=1}^{T}\min(x_t,0)^2}} \label{eq:sortino} \\[4pt]
\text{Max DD} &= \min_{t} \frac{V_t - \max_{s \leq t} V_s}{\max_{s \leq t} V_s} \label{eq:maxdd}
\end{align}

where $\bar{x}$ and $\sigma_x$ are the sample mean and standard deviation of
monthly excess returns and $V_t$ is the portfolio value at month $t$. Drawdowns
are measured on end-of-month values.

%-------------------------------------------------------------
\subsection{Benchmarks}
\label{sec:benchmarks}
%-------------------------------------------------------------

We compare GAT against four benchmarks:
\begin{enumerate}
    \item \textbf{$1/N$ \citep{demiguel2009}}: equal allocation across the 15 risky
          ETFs, rebalanced monthly. \citet{demiguel2009} show this naive rule is
          hard to beat out-of-sample; it is our primary benchmark.
    \item \textbf{60/40}: 60\% SPY + 40\% IEF, rebalanced monthly, the standard
          multi-asset reference.
    \item \textbf{Faber (2007)}: each ETF held equally if its month-end price
          exceeds its 10-month moving average, else the corresponding weight goes
          to BIL \citep{faber2007}, a pure trend peer.
    \item \textbf{S\&P~500 (SPY) buy-and-hold}: the passive equity ceiling.
\end{enumerate}
Statistical comparisons use the paired circular block bootstrap ($b = 6$ months,
$B = 10{,}000$) for Sharpe-ratio differences and Newey--West HAC spanning
regressions for alphas; the same transaction-cost schedule is applied to
GAT and to every benchmark (net versus net).
All Newey--West regressions use $q = 6$ lags throughout the paper (the same
constant across all windows and models), a conservative choice relative to the
rule-of-thumb $\lfloor T^{1/4} \rfloor \approx 4$ for the 60-month test block,
which if anything widens the standard errors and raises the p-values relative to
the rule-of-thumb. Sharpe-difference p-values are robust to block lengths $b \in
\{3,6,12\}$ (see Section~\ref{sec:res_oos}).

%-------------------------------------------------------------
\subsection{Robustness Checks}
\label{sec:robustness}
%-------------------------------------------------------------

Robustness is assessed through:
\begin{enumerate}[label=(\roman*)]
    \item \textbf{Skill-versus-luck attribution.} For each rule we replace its
          informed decision with a random decision drawn from the same feasible
          set at the same frequency (a matched permutation, $N_{\text{sim}} =
          2{,}000$), and test whether the realised Sharpe exceeds the random
          baseline. Only the momentum rule and the minimum-correlation rule beat
          their matched baseline (Section~\ref{sec:results}).
    \item \textbf{Nested incremental spanning-alpha ladder.} A sequence of nested
          models switches on one rule at a time; the Newey--West spanning alpha of
          each model on the previous one isolates that rule's marginal
          contribution.
    \item \textbf{Full-period statistical battery.} Lo (2002) Sharpe standard
          errors and bootstrap CIs, Newey--West HAC tests of the mean excess
          return, Sharpe-difference bootstraps, and CAPM / two-factor / $1/N$
          spanning alphas, on the test block, the post-training sample, and the
          full sample.
    \item \textbf{Transaction costs.} A proportional model on drift-adjusted
          turnover \citep{demiguel2009}, applied identically to all benchmarks,
          with a break-even cost analysis and lower-turnover variants.
    \item \textbf{Construction alternatives.} Holding Rules~1--3 fixed, we replace
          the minimum-correlation rule with a long-only minimum-variance portfolio,
          a risk-parity (inverse-volatility) portfolio, and other constructions over
          the eligible set, to test whether optimised weights dominate the
          equal-weight minimum-correlation triplet (Section~\ref{sec:res_construction}).
    \item \textbf{Deflated Sharpe sensitivity.} The \citet{bailey2014} Deflated
          Sharpe ratio is reported across a range of assumed trial counts, since the
          number of configurations explored is itself a modelling choice.
\end{enumerate}
